% Monaco + WebSocket + 容器内 clangd（与 judge4c 工程一致的概念结构）
\begin{tikzpicture}[
  font=\scriptsize\sffamily,
  box/.style={draw, rounded corners=2pt, align=center, inner sep=5pt,
              minimum width=2.6cm, minimum height=1.05cm},
  arr/.style={-{Stealth[length=2mm]}, thick},
  lbl/.style={font=\tiny, fill=white, inner sep=1pt}
]
  % 主链路：四个节点横向排列，加大间距以容纳标签
  \node[box, fill=blue!8]                          (M)  {浏览器内\\Monaco Editor};
  \node[box, fill=green!10,  right=1.3cm of M]     (C)  {\texttt{MonacoLanguageClient}\\\texttt{vscode-ws-jsonrpc}};
  \node[box, fill=orange!10, right=1.3cm of C]     (P)  {语言服务容器\\jsonrpc-ws-proxy};
  \node[box, fill=purple!8,  right=1.3cm of P]     (D)  {Clangd\\(C/C++ LSP)};

  % 数据库：放在 M 正下方（避开右侧反馈线路径），留足距离
  \node[box, fill=gray!12, below=1.3cm of M] (DB) {数据库\\\texttt{LanguageServerConfig}};

  % 主链路上方：从左到右的请求方向
  \draw[arr] (M) -- node[lbl, above]{LSP JSON-RPC} (C);
  \draw[arr] (C) -- node[lbl, above]{WebSocket}    (P);
  \draw[arr] (P) -- node[lbl, above]{stdio 转发}   (D);

  % 数据库 -> Monaco（SSR 注入端点）
  \draw[arr] (DB) -- node[lbl, right]{SSR 注入端点} (M);

  % 反馈：Clangd -> Monaco（诊断/补全/签名等），从 D 南下绕到 M 南端
  \draw[arr, dashed] (D.south) -- ++(0,-1.0)
                                -| node[lbl, below, pos=0.25]{诊断/补全/签名等}
                                   (M.south);
\end{tikzpicture}
